% =====================================================================
% ejercicio10.tex
% Responsabilidad única: seleccionar y comparar tres modelos bajo un
% costo de negocio asimétrico (L=FP+5FN), y traducir el resultado en
% una recomendación de negocio justificada.
% =====================================================================
\subsection{Ejercicio 10 --- Análisis final y costo de negocio}

Se comparan tres modelos representativos: El de mejor F1/Recall
($k=1$, Euclidiana, estandarizado), el de mejor Accuracy ($k=9$,
Euclidiana, estandarizado), y el modelo base sin estandarizar ($k=5$,
Ejercicio 7).

\begin{table}[H]
\centering
\small
\begin{tabular}{lrrrrrr}
\toprule
Modelo & $k$ & FP & FN & Recall & F1 & Costo $L$ \\
\midrule
Modelo 1 (mejor F1/Recall)          & 1 & 1.293 & 1.575 & 0.2144 & 0.2307 & \textbf{9.168} \\
Modelo 2 (mejor Accuracy)           & 9 &   180 & 1.771 & 0.1167 & 0.1935 & \textbf{9.035} \\
Modelo 3 (base, sin estandarizar)   & 5 &    67 & 1.956 & 0.0244 & 0.0462 & \textbf{9.847} \\
\bottomrule
\end{tabular}
\caption{Comparación final de modelos, con $L = FP + 5FN$.}
\end{table}

\begin{enumerate}

\item \textbf{¿El modelo de mayor Accuracy es también el de menor costo?}
En este caso concreto, \textbf{sí}, ya que el Modelo 2 tiene
simultáneamente el mayor Accuracy ($0.9350$) y el menor costo
($9.035$) de los tres, aunque vale la pena anotarlo con honestidad:
esto \textbf{no contradice} lo demostrado teóricamente en el Ejercicio
5f, donde se construyó un ejemplo específico en el que Accuracy y
costo mínimo sí divergían, sino que simplemente muestra que la
relación entre ambas métricas depende de la forma concreta de los
errores de cada modelo y no es ninguna ley universal, pues aquí el
modelo con más vecinos ($k=9$) resulta lo bastante conservador (con
muy pocos FP) como para también minimizar el costo, algo que podría no
sostenerse con otra función de costo o con otro dataset.

\item \textbf{¿Menor número de falsos negativos?}
El \textbf{Modelo 1} ($k=1$), con $1.575$ FN en vez de los $1.771$ del
Modelo 2 o los $1.956$ del Modelo 3.

\item \textbf{¿Qué modelo recomendaría para identificar clientes de alto riesgo?}
El \textbf{Modelo 1}, a pesar de no tener ni el mayor Accuracy ni el
menor costo $L$, y esto porque si el objetivo explícito es maximizar
la detección de incumplimientos reales (priorizando el Recall), el
Modelo 1 detecta $430$ de los $2.005$ incumplimientos del conjunto de
prueba ($21.44\%$), casi el doble que el Modelo 2 ($234$ de $2.005$,
$11.67\%$), y aunque el costo $L$ tal como está definido pondera los
falsos positivos con peso $1$, esto puede subestimar el valor real de
detectar más incumplimientos si, en la práctica, el costo de no
detectar a un mal pagador termina siendo incluso mayor a $5$ veces el
de una alerta equivocada.

\item \textbf{Justificación final.}
La comparación deja en evidencia una tensión real entre tres criterios
de selección (Accuracy, un costo explícito ponderado, y Recall puro)
que no siempre apuntan al mismo modelo, siendo el Modelo 2 la opción
más conservadora y ``eficiente en costo'' bajo la $L$ definida aquí,
mientras que el Modelo 1 es la más agresiva en detección de riesgo, y
por eso la elección final no es un ejercicio puramente estadístico
sino una decisión de negocio que depende de cuánto le cuesta realmente
a la institución cada tipo de error, ya que si ese costo relativo real
fuera mayor a $5{:}1$, el balance se inclinaría todavía más hacia el
Modelo 1.

\end{enumerate}
